\begin{definition}[Edge-blocked insertion algebra correspondence]%
    \label{def:peps_edgeBlockedInsertionAlgebraIsomorphism_property}
    \lean{TNLean.PEPS.IsEdgeBlockedInsertionAlgebraIsomorphism}
    \leanok
    \uses{def:peps_edgeInsertedCoeff}
    For two PEPS tensors $A$ and $B$ and an edge $e$, this property is
    \begin{align}
        \exists \Phi_e:\MN{D_A(e)}
        \simeq_{\mathrm{alg}}
        \MN{D_B(e)},\qquad
        \forall \sigma,\ \forall X\in\MN{D_A(e)},\quad
        C_A(e,\sigma,X)=C_B(e,\sigma,\Phi_e(X)).
        \notag
    \end{align}
\end{definition}

\begin{theorem}[Edge-blocked insertion algebra isomorphism]%
    \label{thm:peps_edgeBlockedInsertionAlgebraIsomorphism}
    \lean{TNLean.PEPS.isEdgeBlockedInsertionAlgebraIsomorphism}
    \leanok
    \uses{def:peps_edgeBlockedInsertionAlgebraIsomorphism_property,
        thm:peps_edgeBlockedThreeSiteInjective,
        thm:peps_virtualInsertionPhysicalRealization,
        thm:peps_edgeInsertedCoeff_endpointPhysicalRealization,
        thm:peps_edgeInsertedCoeff_injective,
        thm:peps_physical_to_virtual_insertion,
        thm:peps_sameState_edgeBlockedCoeff_eq, def:peps_edgeInsertedCoeff}
    Suppose the two edge-blocked three-site chains associated to $A$ and $B$ at
    the edge $e$ are injective, suppose every virtual bond of $A$ and $B$ has
    positive dimension, and suppose the two PEPS tensors define the same state.
    % Scope restriction: the unrestricted Section 3 statement requires deriving
    % positive bond dimensions from injectivity; see
    % docs/paper-gaps/peps_injective_ft_section3_route.tex.
    Then
    \begin{align}
        \exists \Phi_e:\MN{D_A(e)}
        \simeq_{\mathrm{alg}}
        \MN{D_B(e)},\qquad
        \forall \sigma,\ \forall X\in\MN{D_A(e)},\quad
        C_A(e,\sigma,X)=C_B(e,\sigma,\Phi_e(X)).
        \notag
    \end{align}
    This is the edge-blocked form of the algebra isomorphism $X\mapsto Y$ in
    \cite[Section~3]{Molnar2018NormalPEPS}:
    % Formula: C_A(e,X;sigma)=C_B(e,Y;sigma) for the corresponding three-site
    % blocked chains.  Ink-to-index: X and Y split the distinguished (1,2) bond;
    % the periodic return closes the remaining virtual contraction, and the three
    % upper legs are sigma_1,sigma_2,sigma_3.
    % Boundary signature: both panels have (V,P)=(0,3).  Source verdict: exact
    % reproduction of paper_normal.tex:255--274 and 565--580 in house style.
    \begin{center}
        \tenkzkernel
        \begin{tenkz}[rows={wire,wire}, cols=6, bonds=none, align=1]
            \tn[at=(1,2), name=a1, label pos=s,
                ports={180:virtual, 0:virtual, 90:physical:$\sigma_1$}]{A_1}
            \tn[at=(1,3), skin=ring, name=x, ports={180:virtual, 0:virtual}]{X}
            \tn[at=(1,4), name=a2, label pos=s,
                ports={180:virtual, 0:virtual, 90:physical:$\sigma_2$}]{A_2}
            \tn[at=(1,5), name=a3, label pos=s,
                ports={180:virtual, 0:virtual, 90:physical:$\sigma_3$}]{A_3}
            \tn[at=(2,1), skin=none, name=jw, ports={0:virtual, 90:virtual}]{}
            \tn[at=(2,6), skin=none, name=je, ports={90:virtual, 180:virtual}]{}
            \tnwire{a1.0}{x.180}
            \tnwire{x.0}{a2.180}
            \tnwire{a2.0}{a3.180}
            \tnwire[route=arc]{a3.0}{je.90}
            \tnwire{je.180}{jw.0}
            \tnwire[route=arc]{jw.90}{a1.180}
        \end{tenkz}
        $=$
        \begin{tenkz}[rows={wire,wire}, cols=6, bonds=none, align=1]
            \tn[at=(1,2), name=b1, label pos=s,
                ports={180:virtual, 0:virtual, 90:physical:$\sigma_1$}]{B_1}
            \tn[at=(1,3), skin=ring, name=y, ports={180:virtual, 0:virtual}]{Y}
            \tn[at=(1,4), name=b2, label pos=s,
                ports={180:virtual, 0:virtual, 90:physical:$\sigma_2$}]{B_2}
            \tn[at=(1,5), name=b3, label pos=s,
                ports={180:virtual, 0:virtual, 90:physical:$\sigma_3$}]{B_3}
            \tn[at=(2,1), skin=none, name=jw, ports={0:virtual, 90:virtual}]{}
            \tn[at=(2,6), skin=none, name=je, ports={90:virtual, 180:virtual}]{}
            \tnwire{b1.0}{y.180}
            \tnwire{y.0}{b2.180}
            \tnwire{b2.0}{b3.180}
            \tnwire[route=arc]{b3.0}{je.90}
            \tnwire{je.180}{jw.0}
            \tnwire[route=arc]{jw.90}{b1.180}
        \end{tenkz}
    \end{center}
\end{theorem}
\begin{proof}\leanok
    For a matrix $X\in\MN{D_A(e)}$, realize the insertion of $X$ at the two
    endpoints of the first edge-blocked chain by physical operators $O_1,O_2$.
    Write $C_{\mathrm{mid}}^A$ and $C_{\mathrm{mid}}^B$ for the open middle
    weights of the two edge-blocked chains. For every physical configuration
    $\sigma$, endpoint realization of $X$ and equality of the PEPS states give
    \begin{align}
        C_A(e,\sigma,X)
         & =
        \sum_{\beta}
        O_1(A_u(\beta_u))_{\sigma_u}
        C_{\mathrm{mid}}^A(\sigma;\beta_{\mathrm{L}},\beta_{\mathrm{R}})
                             (A_v(\beta_v))_{\sigma_v}
        \notag \\
         & =
        \sum_{\beta}
        O_1(B_u(\beta_u))_{\sigma_u}
        C_{\mathrm{mid}}^B(\sigma;\beta_{\mathrm{L}},\beta_{\mathrm{R}})
                             (B_v(\beta_v))_{\sigma_v},
        \notag
    \end{align}
    and
    \begin{align}
        C_A(e,\sigma,X)
         & =
        \sum_{\beta}
        (A_u(\beta_u))_{\sigma_u}
        C_{\mathrm{mid}}^A(\sigma;\beta_{\mathrm{L}},\beta_{\mathrm{R}})
        O_2(A_v(\beta_v))_{\sigma_v}
        \notag \\
         & =
        \sum_{\beta}
        (B_u(\beta_u))_{\sigma_u}
        C_{\mathrm{mid}}^B(\sigma;\beta_{\mathrm{L}},\beta_{\mathrm{R}})
        O_2(B_v(\beta_v))_{\sigma_v}.
        \notag
    \end{align}
    Therefore, for every physical configuration $\sigma$,
    \begin{align}
        \sum_{\beta}
        O_1(B_u(\beta_u))_{\sigma_u}
        C_{\mathrm{mid}}^B(\sigma;\beta_{\mathrm{L}},\beta_{\mathrm{R}})
                             (B_v(\beta_v))_{\sigma_v}
         & =
        \sum_{\beta}
        (B_u(\beta_u))_{\sigma_u}
        C_{\mathrm{mid}}^B(\sigma;\beta_{\mathrm{L}},\beta_{\mathrm{R}})
        O_2(B_v(\beta_v))_{\sigma_v}.
        \notag
    \end{align}
    The positive-bond physical-to-virtual recovery applied to the second
    edge-blocked chain gives a matrix $Y\in\MN{D_B(e)}$ such that
    $C_A(e,\sigma,X)=C_B(e,\sigma,Y)$ for every physical configuration
    $\sigma$.
    Thus $X\mapsto Y$ gives the required coefficient correspondence. Denote
    this correspondence by $T$. For every physical configuration $\sigma$,
    \begin{align}
        C_B(e,\sigma,T(X+X'))
         & =C_A(e,\sigma,X+X')
        \notag                               \\
         & =C_A(e,\sigma,X)+C_A(e,\sigma,X')
        \notag                               \\
         & =C_B(e,\sigma,T(X)+T(X')).
        \notag
    \end{align}
    Similarly, $T(cX)=cT(X)$. The multiplicative identity is read from
    the inserted two-endpoint realization:
    \begin{align}
        C_B(e,\sigma,T(XX'))
        =C_A(e,\sigma,XX')
        =C_B(e,\sigma,T(X)T(X')).
        \notag
    \end{align}
    The unit satisfies
    $C_B(e,\sigma,T(1))=C_A(e,\sigma,1)=C_B(e,\sigma,1)$.
    Theorem~\ref{thm:peps_edgeInsertedCoeff_injective} identifies the inserted
    matrices in these coefficient identities. Exchanging $A$ and $B$ gives the
    inverse correspondence, hence a $\C$-algebra isomorphism.
\end{proof}

\begin{theorem}[Matrix-algebra isomorphisms preserve size]%
    \label{thm:matrixAlgEquiv_fin_eq}
    \lean{MPSTensor.matrixAlgEquiv_fin_eq}
    \leanok
    If the full matrix algebras $\MN{D}$ and $\MN{E}$ are isomorphic as
    $\C$-algebras, then $D=E$.
\end{theorem}

\begin{proof}\leanok
    The isomorphism preserves complex vector-space dimension, and
    $\dim_{\C}\MN{D}=D^2$ and $\dim_{\C}\MN{E}=E^2$, so $D=E$.
\end{proof}

\begin{theorem}[Matrix-algebra isomorphisms are inner after identifying size]%
    \label{thm:matrixAlgEquiv_inner_of_fin}
    \lean{MPSTensor.matrixAlgEquiv_inner_of_fin}
    \leanok
    \uses{thm:matrixAlgEquiv_fin_eq, thm:skolem_noether}
    Let $\varphi:\MN{D}\to\MN{E}$ be a $\C$-algebra isomorphism. Then
    $D=E$; let $h$ denote this equality and
    $\iota_h:\MN{D}\to\MN{E}$ the canonical index-relabelling isomorphism.
    There is $X\in\GL_E(\C)$ with
    $\varphi(M)=X\iota_h(M)X^{-1}$ for all $M\in\MN{D}$.
\end{theorem}

\begin{proof}\leanok
    By Theorem~\ref{thm:matrixAlgEquiv_fin_eq}, $D=E$, so $\iota_h$ is defined.
    Then $\varphi\circ\iota_h^{-1}$ is a $\C$-algebra automorphism of
    $\MN{E}$, and Skolem--Noether
    (Theorem~\ref{thm:skolem_noether}) gives $X\in\GL_E(\C)$ with
    $(\varphi\circ\iota_h^{-1})(N)=XNX^{-1}$. Evaluating at
    $N=\iota_h(M)$ gives $\varphi(M)=X\iota_h(M)X^{-1}$.
\end{proof}

\begin{theorem}[Edge gauge from the insertion algebra isomorphism]%
    \label{thm:peps_edgeGaugeFromInsertionAlgebraIsomorphism}
    \lean{TNLean.PEPS.edgeGaugeFromInsertionAlgebraIsomorphism}
    \leanok
    \uses{thm:matrixAlgEquiv_inner_of_fin}
    Suppose the edge insertion correspondence is a $\C$-algebra isomorphism
    between the two full matrix algebras on the chosen bond. Then the two bond
    dimensions on that edge are equal. If $h_e:D_A(e)=D_B(e)$ is this
    equality, and if
    $\phi_{h_e}:\MN{D_A(e)}\to\MN{D_B(e)}$ is the induced algebra
    isomorphism, then, for an invertible edge matrix $Z_e$, the insertion
    correspondence is
    % Formula/source: Papers/1804.04964/paper_normal.tex, lines 254--279
    % and 576--582, identifies the bond insertion map as an algebra
    % isomorphism and writes its inner-conjugation form Y=ZXZ^{-1}.
    % Ink-to-index map: exact algebra is used here because X, phi_h_e(X),
    % and Y are bond endomorphisms, not local tensors with physical legs.
    % Contraction set: the two ordinary matrix products sum the two intermediate
    % bond indices; there are no physical or PEPS lattice contractions.
    % Type verdict: an endomorphism of the B-side edge bond space.
    % Boundary signature: not applicable to the displayed algebra; no tensor
    % boundary is asserted by this formula.
    $Y=Z_e\phi_{h_e}(X)Z_e^{-1}$.
    This is the finite-dimensional matrix-algebra step in
    \cite[Section~3]{Molnar2018NormalPEPS}.
\end{theorem}

\begin{proof}\leanok
    Apply Theorem~\ref{thm:matrixAlgEquiv_inner_of_fin} to the full matrix
    algebras determined by the two edge bond spaces.
\end{proof}

\begin{theorem}[Per-edge gauge family from insertion algebra isomorphisms]%
    \label{thm:peps_exists_edgeGaugeFamily}
    \lean{TNLean.PEPS.exists_edgeGaugeFamily}
    \leanok
    \uses{thm:peps_edgeBlockedThreeSiteInjective,
        thm:peps_edgeBlockedInsertionAlgebraIsomorphism,
        thm:peps_edgeGaugeFromInsertionAlgebraIsomorphism}
    Let $A$ and $B$ be vertex-injective PEPS tensors with the same PEPS state.
    Suppose the corresponding virtual bond dimensions are equal and every
    virtual bond of $A$ has positive dimension. Then there is a family
    $X_e\in\GL(D_A(e))$ such that, for every edge $e$, there is a
    $\C$-algebra isomorphism
    $\Phi_e:\MN{D_A(e)}\to\MN{D_B(e)}$ satisfying
    $C_A(e,\sigma,M)=C_B(e,\sigma,\Phi_e(M))$ for every physical configuration
    $\sigma$ and matrix $M$. Writing
    $\phi_{h_e}:\MN{D_A(e)}\to\MN{D_B(e)}$ for the algebra
    isomorphism induced by the dimension equality on $e$, one has
    $\Phi_e(M)=\phi_{h_e}(X_eMX_e^{-1})$.
\end{theorem}

\begin{proof}\leanok
    Vertex injectivity and positivity give edge-blocked three-site injectivity
    for $A$ and $B$ at each edge.
    Theorem~\ref{thm:peps_edgeBlockedInsertionAlgebraIsomorphism}
    gives the algebra isomorphism $\Phi_e$ preserving inserted coefficients.
    Theorem~\ref{thm:peps_edgeGaugeFromInsertionAlgebraIsomorphism} realizes
    $\Phi_e$ by conjugation, and the bond-dimension equality transports the
    conjugating matrix to the $A$-side bond.
\end{proof}

% Specializations of the per-edge gauge to the translated interior edges of the
% open square lattice, where the gauge comes from a one-edge blocking datum.

\begin{theorem}[Per-edge gauge at a translated horizontal interior edge]%
    \label{thm:peps_exists_regionEdgeGauge_normalSquareHorizontalTranslatedEdge}
    \lean{TNLean.PEPS.exists_regionEdgeGauge_normalSquareHorizontalTranslatedEdge}
    \leanok
    \uses{lem:peps_normalSquareHorizontalTranslatedEdge_blockingData}
    Let $A$ and $B$ be normal PEPS on the open $W\times H$ square lattice that
    satisfy the rectangular-injectivity hypotheses with union closure, have
    positive physical dimension $d>0$, have equal and everywhere-positive virtual
    bond dimensions, and define the same PEPS state. Fix a translated horizontal
    interior frame with two rows and two
    columns of margin, that is $2\le x_0$, $2\le y_0$, $x_0+8\le W$, and
    $y_0+8\le H$, and let $e$ be the distinguished horizontal edge of that frame.
    Then the bond dimensions agree on $e$, $D_A(e)=D_B(e)$, and there are an
    invertible edge matrix $Z\in\GL(D_B(e))$ and a forward map $\mathrm{fwd}$
    given by the conjugation
    $\mathrm{fwd}(M)=Z\phi_{h_e}(M)Z^{-1}$, where
    $\phi_{h_e}:\MN{D_A(e)}\to\MN{D_B(e)}$ is the reindexing induced by the
    bond-dimension equality on $e$. No single-vertex injectivity is used; the
    only injectivity hypotheses are the blocked-region injectivities of the
    interior blocking data and the red and host injectivities of $B$.
\end{theorem}

\begin{proof}\leanok
    The translated horizontal frame supplies, for both $A$ and $B$, the one-edge
    interior blocking datum. Because the red, blue, and complementary blocks are
    coordinate regions, independent of the tensor, the two data share them:
    $R_A=R_B$, $B_A=B_B$, and $C_A=C_B$.
    The distinguished edge $e$ is the unique crossing edge of the red and
    blue blocks. Applying the per-edge gauge construction from a one-edge
    blocking datum to the two blocking data yields the bond-dimension equality on
    $e$ and the conjugating gauge $Z$.
\end{proof}

\begin{theorem}[Per-edge gauge at a translated vertical interior edge]%
    \label{thm:peps_exists_regionEdgeGauge_normalSquareVerticalTranslatedEdge}
    \lean{TNLean.PEPS.exists_regionEdgeGauge_normalSquareVerticalTranslatedEdge}
    \leanok
    The coordinate-swap counterpart of
    Theorem~\ref{thm:peps_exists_regionEdgeGauge_normalSquareHorizontalTranslatedEdge}.
    Under the same hypotheses on $A$ and $B$, for a translated vertical interior
    frame with two rows and two columns of margin, that is $2\le x_0$,
    $2\le y_0$, $x_0+8\le W$, and $y_0+8\le H$, with distinguished vertical
    edge $e$, the bond dimensions agree on $e$, $D_A(e)=D_B(e)$, and there are an
    invertible edge matrix $Z\in\GL(D_B(e))$ and a forward map $\mathrm{fwd}$
    given by the conjugation
    $\mathrm{fwd}(M)=Z\phi_{h_e}(M)Z^{-1}$, with
    $\phi_{h_e}:\MN{D_A(e)}\to\MN{D_B(e)}$ the reindexing induced by the
    bond-dimension equality on $e$.
\end{theorem}

\begin{proof}
    \lean{TNLean.PEPS.Region_map_squareLatticeCoordinateSwap_verticalTranslatedEdgeRed}
    \lean{TNLean.PEPS.Region_map_squareLatticeCoordinateSwap_verticalTranslatedEdgeBlue}
    \leanok
    \uses{thm:peps_exists_regionEdgeGauge_normalSquareHorizontalTranslatedEdge}
    Interchange the two lattice coordinates. This graph isomorphism sends the
    vertical translated frame to the horizontal translated frame, transposes the
    rectangular injectivity hypotheses, and sends the distinguished vertical edge
    to the distinguished horizontal edge. Apply
    Theorem~\ref{thm:peps_exists_regionEdgeGauge_normalSquareHorizontalTranslatedEdge}
    to the transported tensors. Transporting its bond-dimension equality and
    conjugating gauge back along the inverse coordinate swap gives the required
    equality on $e$ and the matrix $Z$.
\end{proof}

\begin{definition}[Factorized local gauge datum]%
    \label{def:peps_hasFactorizedLocalGauge}
    \lean{TNLean.PEPS.HasFactorizedLocalGauge}
    \leanok
    \uses{def:peps_localTensorMap, def:peps_localLeftInverse}
    Fix $A$, $B$, a bond-dimension equality, and a vertex $v$. Let
    $\mathcal T^A_v$ be the local tensor map of $A$, let $L^A_v$ be the
    chosen left inverse of $\mathcal T^A_v$, and let
    $\Pi^A_v=\mathcal T^A_vL^A_v$. The factorized local gauge datum consists
    of the following two identities for local virtual configurations
    $\eta,\eta'$:
    \begin{align}
        \Pi^A_v B_v(\eta,-)
         & =B_v(\eta,-),
        \notag                                 \\
        L^A_vB_v(\eta,-)(\eta')
         & =\prod_{e\ni v}X_e(\eta_e,\eta'_e),
        \notag
    \end{align}
    with $X_e\in\GL(D_A(e))$ on each incident edge. This is the local output
    required from the blocked-middle reduction in
    \cite[Section~3]{Molnar2018NormalPEPS}.
\end{definition}

\begin{definition}[Blocked-middle local gauge formula]%
    \label{def:peps_blockedMiddleGaugeFormula}
    \lean{TNLean.PEPS.BlockedMiddleGaugeFormula}
    \leanok
    The blocked-middle local gauge formula at $v$ is the assertion that there
    are matrices $X_e\in\GL(D_A(e))$, one on each incident edge, such that for
    all local virtual configurations $\eta$ and physical indices $\sigma$,
    \begin{align}
        B_v(\eta,\sigma)
        =
        \sum_{\eta'}
        \left(\prod_{e\ni v}X_e(\eta_e,\eta'_e)\right)
        A_v(\eta',\sigma).
        \notag
    \end{align}
\end{definition}

\begin{theorem}[Blocked-middle formula implies factorized local gauge]%
    \label{thm:peps_hasFactorizedLocalGauge_of_blockedMiddleGaugeFormula}
    \lean{TNLean.PEPS.hasFactorizedLocalGauge_of_blockedMiddleGaugeFormula}
    \leanok
    \uses{def:peps_blockedMiddleGaugeFormula, def:peps_hasFactorizedLocalGauge}
    If the blocked-middle local gauge formula holds at $v$ with matrices
    $X_e$, then the factorized local gauge datum holds at $v$ with the same
    matrices.
    % Open: derive the displayed blocked-middle formula from equality of PEPS
    % states by the edge-centred three-site reduction.
\end{theorem}
\begin{proof}\leanok
    For fixed $\eta$, set
    $c_\eta(\eta')=\prod_{e\ni v}X_e(\eta_e,\eta'_e)$. The displayed
    blocked-middle formula is precisely
    $B_v(\eta,-)=\mathcal T^A_vc_\eta$. Hence
    \begin{align}
        \Pi^A_vB_v(\eta,-)
         & =\Pi^A_v\mathcal T^A_vc_\eta
        =\mathcal T^A_vc_\eta
        =B_v(\eta,-),
        \notag                               \\
        L^A_vB_v(\eta,-)(\eta')
         & =L^A_v\mathcal T^A_vc_\eta(\eta')
        =c_\eta(\eta')
        =\prod_{e\ni v}X_e(\eta_e,\eta'_e).
        \notag
    \end{align}
\end{proof}

\begin{theorem}[Blocked-middle formula gives a local gauge formula]%
    \label{thm:peps_localGauge_exists_of_blockedMiddleGaugeFormula}
    \lean{TNLean.PEPS.localGauge_exists_of_blockedMiddleGaugeFormula}
    \leanok
    \uses{def:peps_blockedMiddleGaugeFormula,
        thm:peps_hasFactorizedLocalGauge_of_blockedMiddleGaugeFormula,
        thm:localGauge_exists}
    Suppose that the edge-centred three-site reduction supplies invertible
    matrices $X_e$ on the incident edges of $v$ such that, for all local virtual
    configurations $\eta$ and physical indices $\sigma$,
    \begin{align}
        B_v(\eta,\sigma)
        =
        \sum_{\eta'}
        \left(\prod_{e\ni v}X_e(\eta_e,\eta'_e)\right)
        A_v(\eta',\sigma).
        \notag
    \end{align}
    Then the same family $X_e$ is a local gauge at $v$:
    % Formula: the displayed component identity above is the local gauge formula.
    % Ink-to-index/type verdict: its boundary is (V,P)=(deg(v),1).  No matching
    % standalone picture occurs in paper_normal.tex, and the old picture hard-coded
    % five virtual legs, so the valence-independent equation is authoritative.
    % Open: derive the blocked-middle formula from SameState
    % by the edge-centred three-site reduction.
\end{theorem}
\begin{proof}\leanok
    By the preceding theorem,
    $L^A_vB_v(\eta,-)(\eta')=\prod_{e\ni v}X_e(\eta_e,\eta'_e)$. Since
    $\Pi^A_vB_v(\eta,-)=B_v(\eta,-)$ and
    $\Pi^A_v=\mathcal T^A_vL^A_v$,
    \begin{align}
        B_v(\eta,\sigma)
         & =\sum_{\eta'}L^A_vB_v(\eta,-)(\eta')A_v(\eta',\sigma)
        \notag                                                   \\
         & =\sum_{\eta'}
        \left(\prod_{e\ni v}X_e(\eta_e,\eta'_e)\right)
        A_v(\eta',\sigma).
        \notag
    \end{align}
\end{proof}
